\section*{Chapter \ref{ch:g6:earth-a-planet} --- The Earth, a Planet of the Solar System}

\begin{solution}{exo:g6:earth-a-planet:1}
\underline{Mercury}, \underline{Venus}, \underline{the Earth},
\underline{Mars}, Jupiter, Saturn, Uranus, Neptune. The \omterm{def:g6:earth-a-planet:planet-types}{rocky planets} are
the first four.
\end{solution}

\begin{solution}{exo:g6:earth-a-planet:2}
``A \omterm{def:g6:earth-a-planet:planet-types}{rocky planet} is small and dense, with a solid surface. A \omterm{def:g6:earth-a-planet:planet-types}{giant planet}
is very large, of low density, made mostly of gases and liquids.''
\end{solution}

\begin{solution}{exo:g6:earth-a-planet:3}
Venus (12\,104~km, against 12\,756~km for the Earth). The largest is
Jupiter (142\,984~km); the smallest is Mercury (4\,879~km).
\end{solution}

\begin{solution}{exo:g6:earth-a-planet:4}
(a) False: Venus has no liquid water; it is far too hot. (b) True
($-65$~$^\circ$C on average, against 15~$^\circ$C). (c) False: the four
planets closest to the Sun are the \omterm{def:g6:earth-a-planet:planet-types}{rocky planets}; the \omterm{def:g6:earth-a-planet:planet-types}{giant planets} are the
four farthest. (d) True: about 0.7~gram per cubic centimetre.
\end{solution}

\begin{solution}{exo:g6:earth-a-planet:5}
$4879 \div 12756 \approx 0.382$, that is 38~\%: Mercury is a little more
than a third as wide as the Earth.
\end{solution}

\begin{solution}{exo:g6:earth-a-planet:6}
$120536 \div 12756 \approx 9.4$: about 9 Earths. Saturn is a \omterm{def:g6:earth-a-planet:planet-types}{giant planet}.
\end{solution}

\begin{solution}{exo:g6:earth-a-planet:7}
Planet~A is a \omterm{def:g6:earth-a-planet:planet-types}{giant planet}: it is very large and its density, 1.1~grams per
cubic centimetre, is low, like that of the \omterm{def:g6:earth-a-planet:planet-types}{giant planets} of the Solar System
(between 0.7 and 1.7). Planet~B is a \omterm{def:g6:earth-a-planet:planet-types}{rocky planet}: it is small and its
density, 5.0, is high, like that of the \omterm{def:g6:earth-a-planet:planet-types}{rocky planets} (between 3.9 and
5.5).
\end{solution}

\begin{solution}{exo:g6:earth-a-planet:8}
(a) $464 - 15 = 449$~degrees hotter. (b) $15 - (-65) = 80$~degrees colder.
\end{solution}

\begin{solution}{exo:g6:earth-a-planet:9}
Venus has a very thick atmosphere of carbon dioxide, which traps the Sun's
heat like a heavy blanket and does not let it escape. Mercury, closer to the
Sun, has almost no air, so it cannot keep its heat.
\end{solution}

\begin{solution}{exo:g6:earth-a-planet:10}
$228 - 150 = 78$~million kilometres farther. $78 \div 150 = 0.52$: Mars's
distance is 52~\% larger than the Earth's.
\end{solution}

\begin{solution}{exo:g6:earth-a-planet:11}
Billions of years ago, Mars had \omterm{def:g1:landscapes-around-us:water}{rivers}, \omterm{def:g1:landscapes-around-us:water}{lakes} and \omterm{def:g5:water-shapes-the-land:flood}{floods}: it was warmer and
wetter. Today its water is frozen, in the ice of the polar caps and under
the ground near the \omterm{def:g3:continents-and-oceans:poles}{poles}. Liquid water cannot last at its surface now
because Mars is very cold ($-65$~$^\circ$C on average) and because its air
is so thin that water quickly freezes or turns into vapour.
\end{solution}

\begin{solution}{exo:g6:earth-a-planet:12}
(a) $8.9 \div 9.8 \approx 0.908$, that is 91~\%. (b) Like a twin: almost
the same size (95~\% of the Earth's diameter), almost the same pull of
gravity, and it is a \omterm{def:g6:earth-a-planet:planet-types}{rocky planet} too. Not like a twin: it is much hotter
(464~$^\circ$C against 15~$^\circ$C), its air is far thicker and made of
carbon dioxide with \omterm{def:g4:clouds-rain-snow:cloud}{clouds} of acid, and it has no liquid water.
\end{solution}

\begin{solution}{pb:g6:earth-a-planet:1}
\textbf{1.} Venus (108~million~km), the Earth (150), Mars (228).

\textbf{2.} $150 - 108 = 42$: Venus is 42~million kilometres closer.
$228 - 150 = 78$: Mars is 78~million kilometres farther.

\textbf{3.} $6792 \div 12756 \approx 0.532$, that is 53~\%: Mars is about
half as wide as the Earth.

\textbf{4.} All three are \omterm{def:g6:earth-a-planet:planet-types}{rocky planets}: their densities (5.2, 5.5 and 3.9
grams per cubic centimetre) are high, and their diameters are small
compared with those of the \omterm{def:g6:earth-a-planet:planet-types}{giant planets}.

\textbf{5.} The thickest air is on Venus (92~bars); the thinnest on Mars
(0.01~bar).

\textbf{6.} $1 \div 0.01 = 100$: the air of Mars presses a hundred times
less than the air of the Earth.

\textbf{7.} Its thick air of carbon dioxide traps the Sun's heat and does
not let it escape, like a heavy blanket.

\textbf{8.} $92 \div 1 = 92$: on the ground of Venus, the air presses
\textbf{92 times} as hard as on the ground of the Earth. A probe must have
very strong walls so as not to be crushed (and must also resist the heat
and the acid \omterm{def:g4:clouds-rain-snow:cloud}{clouds}).

\textbf{9.} The Earth's average temperature, 15~$^\circ$C, is between
0~$^\circ$C and 100~$^\circ$C: water is neither frozen nor boiling, so it
stays liquid.

\textbf{10.} On Venus, at 464~$^\circ$C, far above 100~$^\circ$C, the water
would boil away at once and turn into vapour. On Mars, at $-65$~$^\circ$C,
far below 0~$^\circ$C, it would freeze into ice (and slowly turn into
vapour, because the air is so thin).

\textbf{11.} Long ago, Mars was warmer and wetter than today: liquid water
could flow at its surface. Its \omterm{def:g5:weather-and-climate:climate}{climate} has changed a great deal: it is now
a cold, dry \omterm{def:g5:weather-and-climate:desert}{desert}.

\textbf{12.} For example: ``Dear probes, here on the Earth we have liquid
water, in \omterm{def:g2:water-on-earth:ocean}{oceans}, \omterm{def:g1:landscapes-around-us:water}{rivers} and \omterm{def:g1:landscapes-around-us:water}{lakes}. Our temperatures are mild, 15~$^\circ$C
on average, between the freezing and the boiling of water. We are at the
right distance from the Sun, neither too close nor too far. And our
atmosphere keeps us warm, gives living things the gases they need, and lets
water stay liquid. That is why life is possible here.''
\end{solution}
